% ---------------------------------------------------------------------------
% Preâmbulo da apostila do GHOST.
% Fontes da identidade visual em fontes/.
% Compilar com XeLaTeX.
% ---------------------------------------------------------------------------

\usepackage{fontspec}
\usepackage{unicode-math}
\usepackage{polyglossia}
\setdefaultlanguage{portuguese}

\usepackage[a4paper,top=2.7cm,bottom=2.6cm,left=2.9cm,right=2.9cm,
            headsep=1.0cm,footskip=1.3cm]{geometry}
\usepackage{microtype}
\usepackage{xcolor}
\usepackage{tikz}
\usepackage{pgfplots}
\pgfplotsset{compat=1.18}
\usetikzlibrary{calc,positioning,arrows.meta,decorations.pathreplacing,
                shapes.geometric,backgrounds,fit,patterns}
\usepackage{booktabs}
\usepackage{array}
\usepackage{enumitem}
\usepackage{titlesec}
\usepackage{titletoc}
\usepackage{fancyhdr}
\usepackage{tcolorbox}
\tcbuselibrary{skins,breakable}
\usepackage{caption}
\usepackage[hidelinks]{hyperref}

% ------------------------------------------------------------------ paleta
\definecolor{navy}{HTML}{071B2C}
\definecolor{navyfundo}{HTML}{04121E}
\definecolor{cyan}{HTML}{8DD7DC}
\definecolor{gold}{HTML}{E6B75C}
\definecolor{areia}{HTML}{F6F3EC}
\definecolor{anil}{HTML}{A9AEF0}
% derivadas, para tinta sobre papel claro
\definecolor{tinta}{HTML}{10222F}
\definecolor{cyanescuro}{HTML}{1F6B73}
\definecolor{goldescuro}{HTML}{9A6F1E}
\definecolor{anilescuro}{HTML}{4A4FA8}
\definecolor{grafite}{HTML}{5C6B74}
\definecolor{linha}{HTML}{D6DDD8}
\definecolor{papel}{HTML}{FBFAF6}

% ------------------------------------------------------------------ fontes
\newcommand{\FONTES}{fontes/}
\setmainfont{Newsreader-Regular.ttf}[
  Path            = \FONTES,
  ItalicFont      = Newsreader-Italic.ttf,
  BoldFont        = Newsreader-SemiBold.ttf,
  BoldItalicFont  = Newsreader-SemiBoldItalic.ttf,
  Ligatures       = TeX ]
\setsansfont{Manrope-Regular.ttf}[
  Path            = \FONTES,
  ItalicFont      = Manrope-Regular.ttf,
  BoldFont        = Manrope-SemiBold.ttf,
  BoldItalicFont  = Manrope-SemiBold.ttf,
  Ligatures       = TeX ]
\setmonofont{DMMono-Regular.ttf}[
  Path            = \FONTES,
  ItalicFont      = DMMono-Italic.ttf,
  BoldFont        = DMMono-Medium.ttf,
  BoldItalicFont  = DMMono-MediumItalic.ttf,
  Scale           = 0.86 ]
\newfontfamily\displayfont{Newsreader-Display.ttf}[Path=\FONTES, Ligatures=TeX]
\newfontfamily\manropelight{Manrope-Light.ttf}[Path=\FONTES]
\newfontfamily\manropemed{Manrope-Medium.ttf}[Path=\FONTES]

% Newsreader e DM Mono nao trazem gregas: toda letra grega vem da fonte
% matematica, em modo matematico. Nunca escreva "ν" no texto corrido.
\setmathfont{STIXTwoMath-Regular.otf}

% ------------------------------------------------------------------ titulos
\titleformat{\chapter}[display]
  {\sffamily\bfseries\color{tinta}}
  {\manropelight\fontsize{54}{54}\selectfont\color{cyanescuro}\thechapter}
  {6pt}
  {\fontsize{25}{29}\selectfont\raggedright}
  [{\vspace{4pt}\color{gold}\rule{4.2cm}{2pt}}]
\titlespacing*{\chapter}{0pt}{0pt}{34pt}

\titleformat{\section}{\sffamily\bfseries\large\color{tinta}}{\thesection}{0.7em}{}
\titlespacing*{\section}{0pt}{2.6ex plus .6ex}{1.0ex}
\titleformat{\subsection}{\sffamily\bfseries\normalsize\color{cyanescuro}}{\thesubsection}{0.6em}{}
\titlespacing*{\subsection}{0pt}{2.0ex plus .4ex}{0.6ex}

\setcounter{secnumdepth}{2}
\setcounter{tocdepth}{1}

% ------------------------------------------------------------------ corpo
\setlength{\parskip}{0.55ex}
\setlength{\parindent}{1.1em}
\renewcommand{\arraystretch}{1.24}
\setlist[itemize]{leftmargin=1.3em, itemsep=0.3ex, topsep=0.6ex, label=\color{cyanescuro}\textbullet}
\setlist[enumerate]{leftmargin=1.5em, itemsep=0.4ex, topsep=0.6ex}
\setlist[description]{leftmargin=0pt, font=\sffamily\bfseries\color{tinta}}

% Figuras devem ficar perto do texto que as explica: numa apostila, uma
% figura que migra tres paginas deixa de ser a figura daquele paragrafo.
\renewcommand{\topfraction}{0.92}
\renewcommand{\bottomfraction}{0.85}
\renewcommand{\textfraction}{0.06}
\renewcommand{\floatpagefraction}{0.80}
\setcounter{topnumber}{3}
\setcounter{bottomnumber}{2}
\setcounter{totalnumber}{4}

\captionsetup{font={sf,small}, labelfont={bf,color=cyanescuro}, width=0.92\linewidth}

\pagestyle{fancy}
\fancyhf{}
\renewcommand{\headrulewidth}{0pt}
\fancyhead[L]{\sffamily\scriptsize\color{grafite}\textsc{ghost}}
\fancyhead[R]{\sffamily\scriptsize\color{grafite}\nouppercase{\leftmark}}
\fancyfoot[C]{\sffamily\small\color{grafite}\thepage}
\fancypagestyle{plain}{\fancyhf{}\renewcommand{\headrulewidth}{0pt}%
  \fancyfoot[C]{\sffamily\small\color{grafite}\thepage}}

\hypersetup{colorlinks=true, linkcolor=cyanescuro, urlcolor=cyanescuro,
  pdftitle={GHOST -- uma apostila para comecar},
  pdfsubject={Oscilacao de neutrinos: sabor, materia e decoerencia},
  pdfauthor={Rafael C. R. de Lima}}

% ------------------------------------------------------------------ caixas
\newtcolorbox{nota}[1][]{
  enhanced, colback=papel, colframe=linha, boxrule=0.5pt,
  arc=1.5pt, left=9pt, right=9pt, top=7pt, bottom=7pt,
  borderline west={2.2pt}{-2.2pt}{cyanescuro}, #1}

\newtcolorbox{alerta}[1][]{
  enhanced, colback=papel, colframe=linha, boxrule=0.5pt,
  arc=1.5pt, left=9pt, right=9pt, top=7pt, bottom=7pt,
  borderline west={2.2pt}{-2.2pt}{goldescuro}, #1}

\newtcolorbox{conceito}[1][]{
  enhanced, colback=navy, colframe=navy, boxrule=0pt,
  arc=2pt, left=11pt, right=11pt, top=9pt, bottom=9pt,
  coltext=areia, fontupper=\linespread{1.15}\selectfont, #1}

% rótulo de caixa
\newcommand{\rotulo}[1]{{\sffamily\bfseries\footnotesize\color{cyanescuro}\MakeUppercase{#1}}\par\vspace{2pt}}
\newcommand{\rotuloa}[1]{{\sffamily\bfseries\footnotesize\color{goldescuro}\MakeUppercase{#1}}\par\vspace{2pt}}
\newcommand{\rotuloc}[1]{{\sffamily\bfseries\footnotesize\color{cyan}\MakeUppercase{#1}}\par\vspace{2pt}}

% ------------------------------------------------------------------ atalhos
\newcommand{\nue}{\nu_e}
\newcommand{\numu}{\nu_\mu}
\newcommand{\nutau}{\nu_\tau}
\newcommand{\Losc}{L_{\text{osc}}}
\newcommand{\PH}{P_H}
\newcommand{\dm}[1]{\Delta m^2_{#1}}
\newcommand{\cmd}[1]{{\ttfamily #1}}
\newcommand{\ket}[1]{\left|#1\right\rangle}
\newcommand{\bra}[1]{\left\langle#1\right|}
\newcommand{\braket}[2]{\left\langle#1\middle|#2\right\rangle}

% Espacejamento de letras: no XeLaTeX vem do fontspec, nao do microtype.
\newfontfamily\espacejada{Manrope-SemiBold.ttf}[Path=\FONTES, LetterSpace=22.0]
